% ID: SEL-D0041-OA-E03
% ESTADO: cotejada. Segunda pasada visual y borrador compilado revisados.
% Fuente: Selectividad/matematicas_aplicadas/Ua6MATEMÁTICAS APLICADAS examen.pdf
% SHA-256: b7541cc45d3e96297c29037fe17a562b7daa9cd0ee2a95c2b5e9454ad9b90f14
% Página: 1; opción A; ejercicio original 3.

Un estudio estadístico determina que la noche del 31 de diciembre conduce el 5\% de la población, el 20\% consume alcohol esa noche y el 2\% conduce y consume alcohol.
\begin{enumerate}[label=\alph*)]
\item \textbf{(0.5 puntos)} \textquestiondown{}Son independientes los sucesos ``conducir'' y ``consumir alcohol''?
\item \textbf{(1 punto)} \textquestiondown{}Qué porcentaje de la población no conduce ni consume alcohol esa noche?
\item \textbf{(1 punto)} De las personas que consumen alcohol, \textquestiondown{}qué porcentaje conduce esa noche?
\end{enumerate}
